\documentclass[12pt]{jlreq}
\usepackage{amsmath, amssymb}

\newcommand{\C}{\mathbb{C}}
\newcommand{\R}{\mathbb{R}}
\newcommand{\Z}{\mathbb{Z}}
\newcommand{\N}{\mathbb{N}}
\newcommand{\F}{\mathbb{F}}
\newcommand{\Prob}{\mathbb{P}}
\newcommand{\Exp}{\mathbb{E}}
\newcommand{\dd}{\,d}
\newcommand{\Ker}{\operatorname{Ker}}
\newcommand{\Impart}{\operatorname{Im}}
\newcommand{\Hom}{\operatorname{Hom}}

\begin{document}

閉区間 \([10,20]\) の元 \(t\) に対し，\(\R^3\) 上の \(C^\infty\) 級関数 \(F_t\) を
\[
  F_t(x,y,z)=t(f(x,y))^2+z^2-1,\qquad
  f(x,y)=x^2+y^2(y^2-1)
\]
で定める。このとき，部分集合
\[
  X_t=\{(x,y,z)\in\R^3\mid F_t(x,y,z)=0\}
\]
が \(\R^3\) の \(C^\infty\) 級部分多様体となるような \(t\) の範囲を決定せよ。

\end{document}